% ===== 原文 pp.96–99（PDF p.106–109） =====
% [原文 p.96]
其中 $\delta_h\omega$ 是自旋联络在度规变分下的变化。后一项在度规的任意变分下不消失，因此我们在修正 A-模型的作用量中引入额外的边界项
\[
S_{\partial\Sigma}=\int_{\partial\Sigma}H(\Phi)\,\omega
\]
以确保 $\delta_hS$ 的 Q-正合性。这个边界项不破坏与体运动方程的相容性，因为它不含场的任何导数。它也不破坏作用量的 BRST 不变性，因为 $\chi$ 的边界条件与 $H$ 在 $Y$ 上局部为常数合起来蕴含 $\delta_QH=0$。

% [Tutor Note: p.96 的 $S_{\partial\Sigma}=\int_{\partial\Sigma}H(\Phi)\omega$ 只对边界上有定义，且不含导数，故对 $\Phi$ 的变分给出 $\int_{\partial\Sigma}(\partial_IH)\delta\phi^I\omega$，在 $\chi$ 切于 $Y$、$\dd H|_Y=0$ 的条件下为零。核对 $\delta_h$ 抵消：$\delta_hS_{top}$ 贡献 $-\int_{\partial\Sigma}H\delta_h\omega$，而 $\delta_hS_{\partial\Sigma}=+\int_{\partial\Sigma}H\delta_h\omega$，符号正好相反；若你把 (7.6) 中 $\omega$ 的符号整体反转（等价于把 $\Sigma$ 的定向反转），这个抵消依然成立，但边界项的符号也要跟着翻转。]

\section{等变 Maslov 指标}

让我们重新审视当黎曼面 $\Sigma$ 有非平凡边界时，修正 A-模型路径积分测度的鬼数反常的分析。

\para{带边黎曼面上的丛}

当 $\Sigma$ 有非平凡边界并施加边界条件 (7.9) 时，我们把费米场所生活的丛表述得更精确一些；这里假定子流形 $Y\subset X$ 是拉格朗日的且 $U(1)$ 不变的。

$Y$ 的 $U(1)$ 不变性意味着存在一个相伴的纤维子丛
\[
E_Y=(P|_{\partial\Sigma})\times_{U(1)}Y\quad\subset\quad E|_{\partial\Sigma},
\]
其全空间为
\[
E_Y\equiv\big\{[p,y]\ \big|\ p\in P,\ y\in Y\big\},
\]
投影映射为 $\pi_{E_Y}[p,y]=\pi_E[p,y]=\pi(p)$（记号见附录 C）。玻色场 $\Phi$ 是带边黎曼面 $(\Sigma,\partial\Sigma)$ 上丛对 $(E,E_Y)$ 的截面，即一个映射 $\Phi:\Sigma\to E$，它是投影映射 $\pi_E$ 的右逆，并且满足 $\Phi|_{\partial\Sigma}:\partial\Sigma\to E_Y$。

垂直子丛 $(VE,VE_Y)\subset(TE,TE_Y)$ 分别是 $(\dd\pi_E,\dd\pi_{E_Y})$ 的核。如附录 C 所讨论，$VE$ 从 $X$ 的这些结构继承度规、辛形式与复结构，我们据此分解 $VE\otimes\C=V^{(1,0)}E\oplus V^{(0,1)}E$。

费米场 $\chi$ 是带边黎曼面 $(\Sigma,\partial\Sigma)$ 上一对丛 $(F,F_{\mathbb R})$ 的截面，其中
\begin{gather*}
F\equiv\Phi^*V^{(1,0)}E\\
F_{\mathbb R}\equiv\Phi^*V^{(1,0)}E_Y
\end{gather*}

% [Tutor Note: p.96 的 $(F,F_{\mathbb R})$ 是“实子丛对”的标准记号：$F$ 是 $\Sigma$ 上的全纯向量丛，$F_{\mathbb R}$ 是 $F|_{\partial\Sigma}$ 的一个实子丛，二者构成带边复向量丛对。核对维数：$F$ 的复纤维维数为 $n$，$F_{\mathbb R}$ 的实纤维维数为 $n$，故 p.97 开头写“with real dimensions $(2n,n)$”。当 $Y$ 是拉格朗日且 $U(1)$ 不变时，$F_{\mathbb R}$ 自动是全实的（totally real），这保证 $F_{\mathbb R}\otimes\C=F|_{\partial\Sigma}$。]

% [原文 p.97]
二者的实维数为 $(2n,n)$。事实上，$F_{\mathbb R}$ 是一个全实且拉格朗日的子丛 $F_{\mathbb R}\subset F|_{\partial\Sigma}$（用文献 [22] 的术语），满足
\[
F|_{\partial\Sigma}=F_{\mathbb R}\otimes\C .
\]

\para{全纯向量丛的复双}

与之前一样，鬼数反常将由鬼数为 1 的费米子（$\chi$ 与 $\rho$）的零模数目与鬼数为 $-1$ 的费米子（$\tilde\chi$ 与 $\tilde\rho$）的零模数目不匹配来度量。这又对应于带边黎曼面 $(\Sigma,\partial\Sigma)$ 上全纯向量丛对 $(F,F_{\mathbb R})$ 的 Cauchy-Riemann 算子 $\bar\partial_F$ 的指标（这里的全纯性是就 $\Sigma$ 的内部而言的）。

对丛对计算指标的方法在文献 [21] 中讨论过，其步骤是：先构造 $\Sigma$ 的\emph{复双} $\Sigma_{\C}$，即一个闭黎曼面 $\Sigma_{\C}$，配备反全纯对合 $\tau:\Sigma_{\C}\to\Sigma_{\C}$、二重覆盖 $\pi_\Sigma:\Sigma_{\C}\to\Sigma$，以及一个全纯嵌入 $\Sigma\to\Sigma_{\C}$，使得
\[
\frac{\Sigma_{\C}}{\langle\tau\rangle}=\Sigma,
\qquad\qquad
(\Sigma_{\C})^\tau=\partial\Sigma,
\]
其中 $(\Sigma_{\C})^\tau$ 是对合的不动点集。这个双倍曲面的亏格为
\[
g_{\Sigma_{\C}}=2g_\Sigma+h_\Sigma-1,
\]
其中 $h_\Sigma$ 是边界分支数，$g_\Sigma$ 是洞的数目。（例如，圆盘的复双是配备对径映射 $\tau$、固定赤道的 2-球面。）$\Sigma$ 上的全纯向量丛 $F$ 可以加倍，得到 $\Sigma_{\C}$ 上的全纯向量丛 $F_{\C}$，连同覆盖 $\tau$ 的对合 $\tilde\tau:F_{\C}\to F_{\C}$，并满足
\[
F_{\C}|_\Sigma=F,
\qquad\qquad
(F_{\C})^{\tilde\tau}=F_{\mathbb R}
\]
（见文献 [21] 的定理 3.3.8）。事实上，闭黎曼面 $\Sigma_{\C}$ 上存在一个全纯图册，在其中 $\tilde\tau$ 的作用就是 $\tau$ 作用在底空间上、复共轭作用在纤维 $\C^n$ 上（于是特别地，在这个图册中 $F_{\mathbb R}$ 在 $\partial\Sigma$ 某点上的纤维由实子空间 $\R^n\subset\C^n$ 组成）。

% [Tutor Note: p.97 的亏格公式 $g_{\Sigma_{\C}}=2g_\Sigma+h_\Sigma-1$ 可用两个极端例子核对：圆盘 $g_\Sigma=0,h_\Sigma=1$ 给出 $g=0$（2-球面，与原文括号中的例子一致）；环带 $g_\Sigma=0,h_\Sigma=2$ 给出 $g=1$（环面，复双正是把两条边界粘起来的环面）。另一个通用核对是欧拉示性数：$\chi(\Sigma)=2-2g_\Sigma-h_\Sigma$，而复双的 $\chi(\Sigma_{\C})=2-2g_{\Sigma_{\C}}$，代入上式得 $\chi(\Sigma_{\C})=2\chi(\Sigma)$，这正是一个二重覆盖应有的倍数关系。]
% [Tutor Note: p.97 “(see Theorem 3.3.8 of [21])” 是对文献 [21] 的引用，不是本论文的定理，更不是本章的编号；本文第 7 章没有任何定理类环境（全文只有附录 C 的四个无编号定理、第 9 章的定义 1-7 与命题 1-10、附录 D 的命题）。译文照原样作“见文献 [21] 的定理 3.3.8”，请勿误当作“定理 3.3”。]
% [Tutor Note: p.97 末“闭黎曼面 $F_{\C}$ 上存在一个全纯图册”中，$F_{\C}$ 应为 $\Sigma_{\C}$ 的笔误：全纯图册是底流形的结构，只能长在曲面 $\Sigma_{\C}$ 上，而 $F_{\C}$ 是其上的向量丛（且 $F_{\C}$ 的纤维是 $\C^n$，不是一维复流形）。判断依据：同句后半句说“复共轭作用在纤维 $\C^n$ 上”，可见底空间只能是 $\Sigma_{\C}$。译文按 $\Sigma_{\C}$ 处理。]

% [原文 p.98]
\para{丛对 $(F,F_{\mathbb R})$ 的 Maslov 指标}

作为一个复丛，可以用单个平凡化在 $\Sigma$ 上平凡化 $F$，即 $F\simeq\Sigma\times\C^n$。于是对 $\sigma\in\partial\Sigma$，我们可以用复 $n\times n$ 矩阵 $u(\sigma)\in GL(n,\C)$ 表示 $F_{\mathbb R}$ 在 $\sigma$ 上的纤维（纤维由各列实张成）。假定该平凡化已适配 $F$ 上的度规（从而也适配辛形式），即相应的 Hermite 内积就是 $\C^n$ 上的标准内积；此时刻画 $F_{\mathbb R}\subset F|_{\partial\Sigma}$ 是拉格朗日子丛这一条件相当容易：拉格朗日条件就是
\[
u(\sigma)\in U(n)\subset GL(n,\C)
\]
对每个点 $\sigma\in\partial\Sigma$ 成立。相差 $O(n)$ 元素的两个矩阵表示 $\C^n$ 的同一个实子空间，所以更确切地说，$[u(\sigma)]\in U(n)/O(n)\equiv\Lambda(n)$（$\C^n$ 的拉格朗日 Grassmann 流形）。行列式的平方是从 $\Lambda(n)\to U(1)$ 良定义的映射，它度量该实子空间在 $\C^n$ 中的位置。

于是我们有映射
\[
\partial\Sigma\to\Lambda(n)\to U(1),
\]
而由于 $\partial\Sigma$ 的每个分支是有向圆周，可以对每个分支的上述映射联系一个绕数。对每个边界圆周的这些绕数之和称为丛对 $(F,F_{\mathbb R})$ 的\emph{Maslov 指标}，记作
\[
\mu(F,F_{\mathbb R}).
\]
尽管我们在定义该指标时使用了 $F$ 的一个特定平凡化，但可以证明 [21]，如此得到的整数与平凡化的选取无关。

% [Tutor Note: p.98 用 $\det^2$ 而不是 $\det$ 的理由可直接核对：$u$ 与 $uO$（$O\in O(n)$）表示同一个实子空间，而 $\det(uO)=\det u\cdot\det O$，$\det O=\pm1$；只有平方才在 $O(n)$ 下不变，从而在 $U(n)/O(n)$ 上良定义。另一方面 $U(n)$ 中 $u$ 与 $\ee^{\ii\theta}u$ 给出不同的实子空间（相对于标准复结构差一个整体相位），故 $\det^2$ 恰好记录这个相位——这正是 Maslov 指标（绕数）的几何含义。]
% [Tutor Note: p.98 末“与平凡化无关”可用转移函数的绕数核对：换平凡化 $u\to gu$（$g$ 为 $\partial\Sigma$ 上取值于 $GL(n,\C)$ 的转移函数），$\det^2$ 的绕数改变量是 $\partial\Sigma$ 上 $\det g$ 的绕数的两倍，而这个改变量恰好被 $F$ 的第一陈数在 $\Sigma$ 上的整体贡献吸收，故 $\mu$ 是拓扑不变量。]

\para{指标的计算}

作用在丛对 $(F,F_{\mathbb R})$ 上的 Cauchy-Riemann 算子 $\bar\partial_F$ 的指标，由双倍黎曼面上双倍全纯丛的第一陈数与双倍亏格给出：
\[
\operatorname{ind}\bar\partial_F=c_1(F_{\C})+n(1-g_{\Sigma_{\C}})=c_1(F_{\C})+n(2-2g_\Sigma-h_\Sigma).
\]
整数 $c_1(F_{\C})$ 恰是丛对的 Maslov 指标 $\mu(F,F_{\mathbb R})$。而指标
\[
\mu\big(\Phi^*V^{(1,0)}E,\ \Phi^*V^{(1,0)}E_Y\big)
\]

% [Tutor Note: p.98 的指标公式可作一致性核对：当 $\partial\Sigma=\emptyset$（无边界）时丛对退化为单个丛 $F$，$h_\Sigma=0$，公式给出 $\operatorname{ind}=c_1(F)+n(1-g_\Sigma)$，与 p.93 用 Hirzebruch-Riemann-Roch 得到的结果逐项相同。另外可核对 $n=0$（靶空间为单点）时一切 Maslov 指标为零，与 p.98 的 $\det^2$ 定义（$0\times0$ 矩阵的行列式为 $1$）相容。]

% [原文 p.99]
还可以在 $X$ 是等变 Calabi-Yau 的情形下得到进一步解释：此时 $X$ 上存在一个 $U(1)$ 不变的全纯 $n$-形式，它限制到 $Y$ 上因此在每点 $y\in Y$ 处是体积形式 $\vol_Y$ 的非零复倍数 $c\in\C^\times$。由 $U(1)$ 不变性，$Y$ 上的函数 $c$ 诱导出良定义的函数 $E_Y\to\C^\times$，我们也把它记作 $c$。于是 Maslov 指标就是映射 $c^2\circ\Phi:\partial\Sigma\to\C^\times$ 的绕数。

% [Tutor Note: p.99 用 $c^2$ 而非 $c$，与 p.98 的 $\det^2$ 是同一件事的两种说法：等变 Calabi-Yau 条件给出 $U(1)$ 不变的全纯 $n$-形式 $\Omega_n$，沿 $Y$ 取 $\Omega_n|_Y=c\,\vol_Y$；沿 $E_Y$ 的 $U(1)$ 轨道绕一圈时，$c$ 的相位改变与 $u$ 的行列式的相位改变相同，而平方（两次出现）正是 $\det^2$ 那个在 $O(n)$ 下良定义的量。一致性检验：$n=1$ 时 $c$ 与 p.98 的 $\det u$ 刻画同一个相位，故 p.98 的绕数公式与 p.99 的 $c^2\circ\Phi$ 绕数公式必须给出同一个 $\mu$。]
